%==============================================================
%  lecons/analyse/series_divergentes.tex — Séries divergentes et sommation
%==============================================================

\lecontitle{Séries divergentes}{Analyse réelle — Procédés de sommation de Cesàro et d'Abel}

%=== NOTATION LOCALE =========================================
% (aucune : \sum, \lim sont fournis par amsmath)

%=============================================================
%  § 1 — DÉFINITION ET RÉSULTATS FONDAMENTAUX
%=============================================================
\secnum{navyblue}{1}{Définition et résultats fondamentaux}

\begin{tcolorbox}[thmbox]
  \textbf{Définition (série divergente).}\enspace\index{série divergente}
  Soit $(a_n)_{n\geq 0}$ une suite réelle ou complexe. On lui associe la suite
  des \emph{sommes partielles}
  \[
    S_n \;=\; \sum_{k=0}^{n} a_k .
  \]
  La \emph{série} $\sum a_n$ \emph{converge} si $(S_n)$ admet une limite finie ;
  sinon elle \emph{diverge}. Une série dont le terme général ne tend pas vers $0$
  diverge \emph{grossièrement}.

  \medskip
  \textbf{Définition (sommation de Cesàro).}\enspace\index{sommation de Cesàro}
  On dit que $\sum a_n$ est \emph{sommable au sens de Cesàro}, de somme
  $\sigma\in\mathbb{C}$, si la suite des \emph{moyennes de Cesàro}
  \[
    \sigma_n \;=\; \frac{S_0+S_1+\cdots+S_n}{n+1}
  \]
  converge vers $\sigma$. On note alors $\displaystyle\sum a_n = \sigma\ (C,1)$.

  \medskip
  \textbf{Définition (sommation d'Abel).}\enspace\index{sommation d'Abel}
  On dit que $\sum a_n$ est \emph{sommable au sens d'Abel}, de somme
  $A\in\mathbb{C}$, si la série entière $\sum a_n x^n$ converge pour
  $x\in[0,1[$ et si
  \[
    A \;=\; \lim_{x\to 1^{-}} \sum_{n=0}^{+\infty} a_n x^n .
  \]

  \medskip
  \textbf{Théorème (régularité et hiérarchie).}
  \begin{itemize}
    \item \textit{Condition nécessaire.} Si $\sum a_n$ converge, alors
    $a_n\to 0$. La réciproque est fausse (série harmonique).
    \item \textit{Régularité de Cesàro.} Si $\sum a_n$ converge, de somme $S$,
    alors elle est Cesàro-sommable de même somme $\sigma=S$.
    \item \textit{Régularité d'Abel.} Si $\sum a_n$ converge, de somme $S$,
    alors elle est Abel-sommable de même somme $A=S$ (théorème d'Abel radial).
    \item \textit{Hiérarchie (Frobenius).}\index{théorème de Frobenius} Toute
    série Cesàro-sommable est Abel-sommable, de même somme. La réciproque est
    fausse : la sommation d'Abel est \emph{strictement} plus puissante.
  \end{itemize}
\end{tcolorbox}

\begin{significationintuitive}
Une série peut diverger sans pour autant être « sans valeur » : ses sommes
partielles oscillent, mais leur \emph{moyenne} peut se stabiliser. C'est l'idée
de Cesàro — lisser les oscillations en moyennant. Le procédé d'Abel pondère au
contraire les termes lointains par $x^n\to 0$ et observe la limite radiale. Le
mot clé est la \emph{régularité} : un bon procédé de sommation doit prolonger la
somme usuelle sans la contredire, tout en donnant un sens à des séries qui en
étaient dépourvues.
\end{significationintuitive}

\decsep{}

%=============================================================
%  § 2 — DÉMONSTRATION
%=============================================================
\secnum{navyblue}{2}{Démonstration}

\begin{ideepreuve}
La régularité de Cesàro est le « lemme de Cesàro » : la moyenne d'une suite
convergente converge vers la même limite, par découpage en un bloc fini
négligeable et une queue contrôlée. La régularité d'Abel repose sur une
sommation par parties (transformation d'Abel) reliant $\sum a_n x^n$ aux restes
$S-S_n$. L'inclusion Cesàro~$\Rightarrow$~Abel s'obtient en exprimant la série
entière à l'aide des moyennes $\sigma_n$.
\end{ideepreuve}

\vspace{10pt}
\rubriq{Preuve}

\begin{mdframed}[style=proofbar]
  \textit{Étape 1 — Condition nécessaire.}
  Si $S_n\to S$, alors $a_n=S_n-S_{n-1}\to S-S=0$. La série harmonique
  $\sum 1/n$ montre que la réciproque est fausse : $1/n\to 0$ mais
  $S_{2n}-S_n\geq n\cdot\tfrac{1}{2n}=\tfrac12$, donc $(S_n)$ n'est pas de
  Cauchy.

  \medskip
  \textit{Étape 2 — Régularité de Cesàro.}
  Supposons $S_n\to S$. Quitte à remplacer $a_0$ par $a_0-S$, on se ramène à
  $S=0$. Soit $\varepsilon>0$ ; il existe $N$ tel que $|S_n|\leq\varepsilon$
  pour $n\geq N$. Pour $n>N$,
  \[
    |\sigma_n|
    \;\leq\; \frac{1}{n+1}\sum_{k=0}^{N}|S_k|
      \;+\; \frac{1}{n+1}\sum_{k=N+1}^{n}|S_k|
    \;\leq\; \frac{C_N}{n+1} \;+\; \varepsilon ,
  \]
  où $C_N=\sum_{k=0}^{N}|S_k|$ est fixe. En faisant $n\to+\infty$ on obtient
  $\limsup|\sigma_n|\leq\varepsilon$, ceci pour tout $\varepsilon>0$, donc
  $\sigma_n\to 0=S$.

  \medskip
  \textit{Étape 3 — Régularité d'Abel.}
  Supposons $S_n\to S$. La suite $(S_n)$ est bornée, donc $\sum a_n x^n$ a un
  rayon de convergence $\geq 1$. Pour $x\in[0,1[$, la transformation d'Abel
  $a_n=S_n-S_{n-1}$ (avec $S_{-1}=0$) donne
  \[
    \sum_{n=0}^{+\infty} a_n x^n
    \;=\; (1-x)\sum_{n=0}^{+\infty} S_n x^n .
  \]
  Comme $(1-x)\sum_{n\geq 0} x^n = 1$, on écrit
  \[
    \sum_{n=0}^{+\infty} a_n x^n - S
    \;=\; (1-x)\sum_{n=0}^{+\infty} (S_n-S)\,x^n .
  \]
  Soit $\varepsilon>0$ et $N$ tel que $|S_n-S|\leq\varepsilon$ pour $n\geq N$.
  En séparant les deux blocs,
  \[
    \Bigl|\sum a_n x^n - S\Bigr|
    \;\leq\; (1-x)\sum_{n=0}^{N}|S_n-S|
      \;+\; \varepsilon\,(1-x)\sum_{n>N} x^n
    \;\leq\; (1-x)\,C_N' + \varepsilon ,
  \]
  avec $C_N'=\sum_{n=0}^{N}|S_n-S|$. En faisant $x\to 1^{-}$, le premier terme
  tend vers $0$, d'où $\limsup_{x\to 1^-}|\sum a_n x^n-S|\leq\varepsilon$. Donc
  la limite radiale vaut $S$.

  \medskip
  \textit{Étape 4 — Cesàro $\Rightarrow$ Abel.}
  Supposons $\sigma_n\to\sigma$. Puisque $(\sigma_n)$ converge, on a
  $S_n=(n+1)\sigma_n-n\sigma_{n-1}=\mathcal{O}(n)$, puis
  $a_n=S_n-S_{n-1}=\mathcal{O}(n)$ : les séries entières $\sum a_n x^n$ et
  $\sum S_n x^n$ ont donc un rayon de convergence $\geq 1$, et la
  transformation d'Abel de l'étape~3 reste valable sur $[0,1[$, soit
  $\sum_{n\geq 0}S_n x^n
  =\tfrac{1}{1-x}\sum_{n\geq 0}a_n x^n$. De même, en appliquant la
  transformation d'Abel à la suite $(S_n)$ dont $(n+1)\sigma_n=S_0+\cdots+S_n$
  est la suite des sommes partielles,
  $\sum_{n\geq 0}S_n x^n=(1-x)\sum_{n\geq 0}(n+1)\sigma_n x^n$. En combinant,
  \[
    \sum_{n=0}^{+\infty} a_n x^n
    \;=\; (1-x)^{2}\sum_{n=0}^{+\infty} (n+1)\,\sigma_n\, x^n .
  \]
  Puisque $(1-x)^2\sum_{n\geq 0}(n+1)x^n=1$, le même argument de découpage qu'à
  l'étape~3 (appliqué à la suite convergente $(\sigma_n)$, avec le poids
  $(1-x)^2(n+1)x^n$ de somme $1$) donne
  $\lim_{x\to1^-}\sum a_n x^n=\sigma$. Ainsi la série est Abel-sommable de
  somme $\sigma$.
  \hfill$\blacksquare$
\end{mdframed}

\decsep{}

%=============================================================
%  § 3 — EXEMPLES, CONTRE-EXEMPLES ET EXERCICES
%=============================================================
\secnum{exgreen}{3}{Exemples, contre-exemples et exercices}

\rubriq{Exemples}

\begin{mdframed}[style=exbar]
  \textbf{1. Série de Grandi.}\enspace\index{série de Grandi}%
  Pour $a_n=(-1)^n$, les sommes partielles valent $S_n=1,0,1,0,\dots$ La série
  diverge, mais les moyennes de Cesàro $\sigma_n\to\tfrac12$ et la série
  d'Abel $\sum(-1)^n x^n=\tfrac{1}{1+x}\to\tfrac12$. Donc
  \[
    1-1+1-1+\cdots \;=\; \tfrac12 \quad (C,1)\ \text{et}\ (A).
  \]

  \smallskip\noindent
  \textbf{2. Une série Abel-sommable mais non Cesàro-sommable.}\enspace%
  Pour $a_n=(-1)^n(n+1)$, on a $\sum a_n x^n=\tfrac{1}{(1+x)^2}\to\tfrac14$,
  donc $1-2+3-4+\cdots=\tfrac14\ (A)$. En revanche les moyennes $\sigma_n$
  divergent : la série n'est pas Cesàro-sommable. Ceci illustre la stricte
  hiérarchie Cesàro $\subsetneq$ Abel.

  \smallskip\noindent
  \textbf{3. Théorème de Frobenius en action.}\enspace%
  Toute série convergente est à la fois Cesàro- et Abel-sommable, de même
  somme : les procédés \emph{prolongent} la sommation usuelle. Par exemple
  $\sum(-1)^n/(n+1)=\ln 2$ converge, donc vaut $\ln 2$ aussi au sens d'Abel,
  ce qui redonne $\ln 2=\lim_{x\to1^-}\ln(1+x)$.
\end{mdframed}

\vspace{6pt}
\rubriq{Contre-exemples et pièges}

\begin{mdframed}[style=cexbar]
  \textbf{Divergence grossière non sommable.}\enspace%
  La série $1+2+3+\cdots$ ($a_n=n$) n'est ni Cesàro- ni Abel-sommable :
  $\sum n\,x^n=\tfrac{x}{(1-x)^2}\to+\infty$. L'écriture « $=-1/12$ » relève du
  prolongement de la fonction zêta, \emph{pas} d'un procédé de sommation
  régulier.\index{prolongement analytique}

  \smallskip\noindent
  \textbf{Linéarité, pas réarrangement.}\enspace%
  Les sommes $(C,1)$ et $(A)$ sont linéaires et stables par décalage, mais
  \emph{pas} invariantes par permutation des termes. Réarranger
  $1-1+1-1+\cdots$ peut changer la moyenne de Cesàro.

  \smallskip\noindent
  \textbf{La réciproque d'Abel est fausse (Tauber).}\enspace%
  Abel-sommable n'implique pas convergente : $\sum(-1)^n$ en est l'exemple. Il
  faut une hypothèse \emph{taubérienne} (par exemple $a_n=o(1/n)$, théorème de
  Tauber) pour récupérer la convergence à partir de la sommabilité d'Abel.%
  \index{théorème taubérien}
\end{mdframed}

\vspace{6pt}
\exolabel

\begin{mdframed}[style=exobar]
  \begin{enumerate}
    \item Montrer que la série harmonique $\sum 1/n$ diverge, mais que
          $\sum (-1)^{n-1}/n$ converge. La première est-elle Cesàro-sommable ?

    \item (Lemme de Cesàro) Démontrer en détail que si $u_n\to\ell$, alors
          $\tfrac{1}{n+1}\sum_{k=0}^{n}u_k\to\ell$. En déduire la régularité de
          la sommation de Cesàro.

    \item Calculer les moyennes de Cesàro de la série de Grandi et retrouver
          $\sigma_n\to\tfrac12$. \textit{Indication :} distinguer $n$ pair et
          impair.

    \item Établir l'identité $\sum_{n\geq0}a_n x^n=(1-x)\sum_{n\geq0}S_n x^n$
          pour $x\in[0,1[$ (transformation d'Abel) et en déduire le théorème
          d'Abel radial.

    \item Vérifier que $\sum(-1)^n(n+1)x^n=(1+x)^{-2}$, puis que cette série
          n'est pas Cesàro-sommable bien qu'Abel-sommable de somme $\tfrac14$.

    \item (Ouverture — théorème taubérien de Tauber) Soit $\sum a_n$
          Abel-sommable de somme $A$, avec $n a_n\to 0$. Montrer que $\sum a_n$
          converge et que sa somme vaut $A$. \textit{Indication :} encadrer
          $S_n-\sum a_k x^k$ pour $x=1-\tfrac1n$.
  \end{enumerate}
\end{mdframed}

\leconfooter{E.\ Cesàro (1890) ; N.\,H.\ Abel (1826) ; A.\ Tauber (1897)}
